\section{Diseño de gabinetes de módulos}

Para la construcción de los gabinetes de cada módulo y del gabinete general del prototipo, se opta por la impresión 3D, ya que esta tecnología permite fabricarlos a medida con las especificaciones requeridas, además de ofrecer costos accesibles y una rápida producción. Los diseños se realizan en FreeCAD \parencite{FreeCAD}, seleccionado por ser un software libre, fácil de usar y altamente versátil.

\subsection{Gabinete general}

El gabinete principal, encargado de controlar todo el sistema, debe tener la capacidad suficiente para alojar en su interior la fuente de tensión de \SI{12}{\volt} de corriente continua, el transformador de \SI{48}{\volt} de corriente alterna y la PCB general de control y distribución de energía. Además, en la superficie debe quedar espacio para ubicar el gabinete individual en prueba, junto con la pantalla interactiva y el \textit{encoder} rotativo.

En base a estos requerimientos, se diseña un gabinete compuesto por dos partes: una inferior (base), destinada a la fijación de los elementos de alimentación y control, y una superior, en la cual se instala el display de \SI{2.8}{''} con el \textit{encoder}, además del espacio para apoyar el gabinete individual.  

Asimismo, se diseña un soporte en ángulo para el display, que cumple la doble función de mejorar su visualización y proteger la PCB del mismo.

\begin{figure}[H]
    \centering
    \includegraphics[width=0.4\textwidth]{3d/Angulo_pantalla.png}
    \caption{Soporte en ángulo para el display.}
    \label{fig:3d:soporte_pantalla}
\end{figure} 

Las dimensiones finales del gabinete son de $20\times15 \, \unit{\centi\meter}$, con sus piezas unidas mediante tornillos. En la base se realizan todas las perforaciones necesarias para fijar la PCB, la fuente de corriente continua y el transformador de corriente alterna con tornillos Allen de cabeza fresada. Finalmente, en el lateral posterior se incorpora una rejilla de ventilación y dos aberturas para instalar un conector macho \textit{interlock} y un interruptor de encendido. En el lateral izquierdo se deja una ventana destinada al paso del cable plano que conecta el gabinete general a los módulos.

\begin{figure}[H]
    \centering
    \begin{subfigure}[b]{\textwidth}
        \centering
        \includegraphics[width=0.6\textwidth]{3d/Gabinete_general_base.png}
        \caption{Base de gabinete general.}
        \label{fig:3d:base_gabinete_general}
    \end{subfigure}
    \begin{subfigure}[b]{\textwidth}
        \centering
        \includegraphics[width=0.6\textwidth]{3d/Gabinete_general_superior.png}
        \caption{Tapa superior de gabinete general.}
        \label{fig:3d:superior_gabinete_general}
    \end{subfigure}
    \caption{Piezas que componen el gabinete general.} 
    \label{fig:3d:gabinete_general}
\end{figure}

\subsection{Gabinete individual}

El gabinete destinado a cada módulo se concibe inicialmente como una caja cuadrada en la cual pudiera fijarse una placa de circuito impreso (PCB) de $10\times10 \, \unit{\centi\meter}$, separada un centímetro del fondo mediante cuatro separadores huecos que permiten el paso de tornillos. En un primer diseño, la caja no contemplaba tapa, lo que implicaba un contacto directo del usuario con la PCB. Para evitarlo, se incorpora una tapa de acrílico transparente que cubre la totalidad de la superficie superior del gabinete. De esta forma se impide el contacto directo con la electrónica, pero al mismo tiempo se mantiene la visibilidad de la placa. No obstante, esta solución dificulta la realización de mediciones directas, por lo que se añadieron perforaciones en el acrílico, de tamaño suficiente para permitir el ingreso de las puntas de un multímetro u osciloscopio, sin comprometer la seguridad del usuario.

Posteriormente se considera que los módulos solo requieren de una PCB de $10\times5 \, \unit{\centi\meter}$. Sin embargo, se decide mantener el diseño original para conservar una línea de universalidad que brinde libertad al usuario en el desarrollo de distintos módulos. Para ello, se añaden dos separadores adicionales que permiten fijar también placas de $10\times5 \, \unit{\centi\meter}$ por sus cuatro esquinas.

La última consideración es la conexión de los módulos. Se busca una solución que evitara tener que desmontar la PCB para realizar las conexiones. Esto se resuelve diseñando el gabinete con una ventana lateral del tamaño exacto del conector, de modo que la PCB pueda alojar el conector en su borde y conectarse directamente desde el exterior.

\begin{figure}[H]
    \centering
    \includegraphics[width=0.6\textwidth]{3d/Gabinete_chico.png}
    \caption{Gabinete individual.}
    \label{fig:3d:gabinete_individual}
\end{figure}

\subsubsection{Conector PCIe x4}

El conector PCIe x4 se monta sobre una PCB desde la cual se desprenden los cables correspondientes. Para su protección, se diseña una caja negra en la que, en su cara frontal, queda accesible el conector; mientras que en la cara posterior se coloca una tapa con una abertura lo suficientemente amplia para permitir el paso de los cables.

\begin{figure}[H]
    \centering
    \includegraphics[width=0.5\textwidth]{3d/Carcasa_conector.png}
    \caption{Carcasa para conector PCIe x4.}
    \label{fig:3d:carcasa_conector}
\end{figure}

